% Terminos (glosario general)
\newglossaryentry{helada}{
  name={helada},
  description={Fenómeno meteorológico en el que la temperatura del aire desciende a 0\,°C o menos, dañando los tejidos vegetales}
}

\newglossaryentry{seriestemporales}{
  name={series temporales},
  description={Secuencia de datos ordenados en el tiempo, usada para detectar patrones y hacer predicciones}
}

\newglossaryentry{desbalance}{
  name={desbalance de clases},
  description={Situación en que una clase (helada) es mucho menos frecuente que la otra (no helada), lo que dificulta el entrenamiento de los modelos}
}

\newglossaryentry{sobreajuste}{
  name={sobreajuste},
  description={Situación en la que un modelo memoriza los datos de entrenamiento y falla al generalizar a datos nuevos}
}

\newglossaryentry{persistencia}{
  name={modelo de persistencia},
  description={Modelo base que predice que las condiciones de mañana serán iguales a las de hoy}
}

% Abreviaturas (siglas)
\newacronym{senamhi}{SENAMHI}{Servicio Nacional de Meteorología e Hidrología del Perú}
\newacronym{ml}{ML}{Aprendizaje Automático (\textit{Machine Learning})}
\newacronym{lstm}{LSTM}{Long Short-Term Memory}
\newacronym{xgb}{XGBoost}{Extreme Gradient Boosting}
\newacronym{rf}{RF}{Random Forest}
\newacronym{aucpr}{AUC-PR}{Área bajo la curva Precisión-Recall}
\newacronym{mcc}{MCC}{Coeficiente de correlación de Matthews}

% Simbolos
\newglossaryentry{simboloC}{
  name={$^\circ$C},
  description={Grados Celsius, unidad de temperatura},
  type=symbols
}

\newglossaryentry{simboloTmin}{
  name={$T_{\min}$},
  description={Temperatura mínima diaria del aire},
  type=symbols
}